\section{Limitations and Remaining Work}
\label{sec:limitations}

We state the current status of each limitation explicitly, including the two that
were previously hard submission gates.

\textbf{The evaluation cohort had to be corrected for a subject leak, and this
matters.} The de Chazal DS1/DS2 split is record-disjoint but not
subject-disjoint: records 201 (DS1) and 202 (DS2) are the same person, per
record 202's own WFDB header. Any result computed on the full 22-record DS2 set
includes one patient the source model was trained on. We audited all 48 records,
found this to be the only such pair, and excluded record 202. The correction is
not cosmetic --- it removed the statistical significance of the A4--A3 contrast
(Holm $p$ $0.047 \to 0.075$), leaving only the two comparisons against the frozen
model surviving correction. Published results on this split that do not exclude
202 are not strictly patient-disjoint and may be optimistically biased by the
same mechanism, and we recommend the check
be run as standard.

\begin{sloppypar}
\textbf{Robustness to source-training objective (RQ5) --- partially addressed.}
Five architecture-identical DS1 source variants (cross-entropy, effective-number
weighting, focal $\gamma\in\{1,2,3\}$) are trained with five seeds and reach
S sensitivity $0.78$--$0.88$ on their own DS1 test split, a figure that motivated
the concern. The adaptation gain persists under two alternative
class-reweighted source-training objectives (Table~\ref{tab:strong_adapt}).
This supports robustness to the choice of source \emph{objective}, but does
\textbf{not} establish robustness to a demonstrably stronger unseen-patient
source model. Three gaps remain.
Adaptation was re-run from two of the five trained
variants, so the robustness claim covers two sources rather than five. The three
RR-placement and CNN variants (RR-conditioned token embedding, RR token before
attention, and a lightweight inter-patient CNN) require new model graphs and
were never trained. The encoder-minus-head margins from the reweighted
checkpoints, while nonnegative, are small and carry no reported uncertainty, so
RQ5 does not independently support an encoder-over-head ordering.
 These
checkpoints are also not \emph{stronger} sources --- on intra-patient DS1 S
sensitivity they are below the cross-entropy baseline --- so RQ5 tests
robustness to the source objective, not adaptation from a stronger start.
\end{sloppypar}

\textbf{Cross-database transfer (RQ6) --- closed at subject level, with a
cohort caveat.} INCART recordings were grouped by the subject identifier in the
WFDB header and the saved cells re-aggregated to one score per subject, so the
external analysis no longer treats repeated recordings from one person as
independent. The gain survives the correction: $+0.160$ \macrof{} for encoder
LoRA over 6 subjects with a bootstrap interval excluding zero and 6/6 subjects
improved, against $+0.226$ over 46 subjects on SVDB. Two limits remain. With
$n = 6$ the INCART result is descriptive rather than inferential, and the INCART
cohort is a tractability-capped convenience sample --- only the first 13 of 75
records were attempted --- so it is not a principled sample of the database. Both
external cohorts now have itemized inclusion and exclusion manifests
(Table~\ref{tab:external_manifest}) in which every configured recording carries
exactly one deterministic disposition. Source-only cross-database generalization
remains weak; it is the target-informed adaptation that is positive.

\textbf{The budget sweep is partial, and named accordingly (RQ4).} Of 140 planned budget $\times$
optimizer $\times$ rank $\times$ replay-location configurations, 49 are measured,
5 are analytically infeasible, and 86 are unrun (Table~\ref{tab:exp_status}). The measured surface is complete
only for Adam at ranks 1--2, and only on a reduced 12-record, 3-seed panel. This
supports only the narrow claim we make---that \emph{within the measured
12-record, three-seed Adam rank-1/rank-2 panel, increasing the budget from 8 to
64\,KiB produced no statistically detectable benefit}, tested by paired contrasts
and TOST equivalence rather than by interval overlap---and supports nothing about
optimizer choice, higher ranks, or budgets outside that panel. Numbers from
this panel are never compared against primary-cohort arms. Following the
completed scope, we describe E10 as an \emph{Adam budget sweep for rank-1 and
rank-2 adaptation}, not as the full optimizer grid it was originally designed
to be.

\textbf{Remaining limitations.} (i)~The primary-cohort direct A4/A5-vs-A1
comparison remains non-significant ($n=21$), so a definitive ``plasticity beats
volume'' claim is not licensed; a larger patient cohort is the most plausible
route to settling it, and the cohort is now one record smaller than the
protocol's nominal size. (ii)~Persistent-state totals are packed and asserted per cell, and the whole grid
was replicated against a 1\,KiB firmware implementation reserve with no change in
conclusions (Section~\ref{sec:reserve}); a byte-exact MCU port would still add
allocator overhead not modeled here. (iii)~Peak SRAM, update latency, update
energy, and Flash write volume are all unmeasured, and no hardware training loop
was executed. \textbf{Training compute, update latency, cycles, and energy are
unmeasured, and no numerical compute ratio is reported}, because the
encoder-LoRA backward cost requires layer-level profiling or a complete
reverse-mode derivation that this study does not provide
(Appendix~\ref{sec:embedded}). (iv)~The regularization family is represented
here by \emph{one inexpensive, memory-matched $B$-factor anchor}
(Section~\ref{sec:reg}), chosen because full-model FP32 EWC costs
$53{,}144$\,B and is infeasible at this budget. Note it anchors the $B$ factor,
not the functional update $BA$. Restricted EWC, distillation from the frozen
model, output-consistency regularization, and functional-delta ($\|AB\|_F^2$)
regularization all remain untested, so this is a comparison against one
convenient member of the family and not against the family. (v)~MIT--BIH is small and old, and Q/paced beats are unsupported in the
primary cohort. (vi)~The reproducibility audit (Table~\ref{tab:audit}) returns 11 PASS, 0 FAIL,
and one WARN: configuration values are fixed in a versioned config file, but no automated
provenance proves the hyperparameters were chosen without access to final test
results. We report the WARN rather than suppressing it; the mitigation is that
hyperparameters are shared across all arms, so any residual selection effect
applies to every arm equally and cannot by itself produce the between-arm
ordering. (vii)~The method assumes supervised labels for an early patient segment
and has had no prospective clinical evaluation.

\textbf{Framing rule.} Under a fixed analytical adaptation-state budget,
patient-specific adaptation of the tiny ECG transformer exhibits a
\emph{replay--plasticity allocation frontier}. Encoder low-rank adaptation with
limited raw replay significantly improves over the frozen model, produces
stronger pooled minority-class performance, and yields more threshold-free
positive patient changes than maximum-volume post-pool replay. However, the
direct average difference from maximum-volume head-only replay is
\textbf{not statistically significant}. The evidence therefore supports co-design
and several competitive operating points rather than a universal winner, and no
sentence in this paper should be read as establishing superiority over A1.
